\section{Introduction}\label{sec:introduction}

Let $G$ be a simply connected complex semisimple group, and write $V(\lambda)$
for its irreducible representation of dominant integral highest weight
$\lambda$. The saturation problem asks whether occurrence of a tensor
invariant at dilated highest weights forces occurrence at the original
weights. A necessary congruence is that the sum of the original weights
belong to the root lattice: otherwise the center acts nontrivially on the
tensor product. Write $Q$ for the root lattice. The question is therefore
whether the semigroup of triples with a tensor invariant contains every
triple of dominant integral weights in its rational cone whose sum
belongs to $Q$.

This arithmetic question is closely connected with the geometry of
eigenvalues. For compact groups, the corresponding real cone describes
which three coadjoint orbits contain representatives with sum zero.
Klyachko connected the Hermitian eigenvalue problem with tensor
occurrence after dilation \cite{Kly98}. The relation between moment
polytopes and occurrence after integral
dilation is part of the general branching framework of Berenstein and
Sjamaar \cite[Theorems~3.2.1--3.2.2]{BS00}. Saturation asks whether this
real feasibility, together with the central congruence, already detects
an invariant at the original weights.

Knutson and Tao's honeycomb theorem answered this question in type $A$
\cite{KT99}. Derksen and Weyman gave a quiver proof through saturation
of semi-invariant weights \cite{DW00}, and Belkale gave a geometric
proof \cite{Bel06}. Kapovich and Millson formulated the analogous root-lattice
assertion for simply laced groups
\cite[Conjecture~1.4(1)]{KM06}; see also
\cite[Conjecture~1.5]{KKM09}. We prove its type-$D$ case.

\begin{theorem}[Type-$D$ saturation]\label{thm:main}
Let $G=\Spin(2n)$, $n\ge2$, and let $\lambda,\mu,\nu$ be dominant
integral weights such that $\lambda+\mu+\nu\in Q(D_n)$. Then, for every
positive integer $N$,
\[
 \bigl(V(\lambda)\otimes V(\mu)\otimes V(\nu)\bigr)^G\ne0
 \quad\Longleftrightarrow\quad
 \bigl(V(N\lambda)\otimes V(N\mu)\otimes V(N\nu)\bigr)^G\ne0.
\]
\end{theorem}

The theorem includes the half-integral spin weights. We use the standard
convention $D_2=A_1\times A_1$; $D_3$ is isomorphic to $A_3$. If $n=1$
is included, $G$ is a torus and the assertion is just the weight-sum
equation. The stretching parameter throughout is a positive integer.

Kapovich, Kumar, and Millson proved saturation
for $\Spin(8)$ \cite[Theorem~5.3]{KKM09}, and Kiers proved it for $\Spin(10)$ and
$\Spin(12)$ \cite[Theorem~1.3]{Kiers19}. These results include the
spin representations and establish the first nontrivial ranks beyond
the type-$A$ identifications. General saturation-factor results give
another important approach. The building and path methods of
Kapovich and Millson give a factor of four in type $D$
\cite[Theorem~1.1]{KM08}. Belkale and Kumar obtained factor two for
symplectic and odd special orthogonal groups \cite{BK10}. Sam's
symmetric-quiver methods give factor two uniformly for special
orthogonal and symplectic groups and, by descent of doubled weights,
factor four for spin groups
\cite[Theorem~1.1 and Corollary~1.2]{Sam12}. These factors do not give
the scale-one conclusion under the root-lattice hypothesis.
Hong and Shen subsequently obtained factor two for odd spin groups
under the root-lattice congruence \cite[Corollary~1.3]{HS15}.
Recent progress also includes Min's proof of saturation for
Newell--Littlewood numbers for triples of partitions of even total
size \cite[Theorem~1.2]{Min24}. These numbers describe stable
ordinary-weight tensor multiplicities for the classical groups;
that range excludes general spin weights and the full unstable
range covered here.

Our proof uses paths, but separates their geometric and arithmetic
roles from the construction of representations. Littelmann's path
model encodes tensor products through root operators and integral
saturated chains \cite{Lit94,Lit95}. Kapovich and Millson relate such
paths to geodesics in buildings \cite{KM08}, while Kapovich, Leeb, and
Millson connect weighted configurations with polygons in symmetric
spaces \cite{KLM09}. We use this geometric viewpoint to obtain real
paths, then solve the type-$D$ lattice problem by deforming their
switches. The resulting integral schedules have a weaker contract than
Lakshmibai--Seshadri paths. A separate construction with exterior and
spin columns supplies tensor occurrence.

We describe the main constructions. Put $\xi=\nu^*$, so that the
desired invariant is equivalent to an intertwiner
$V(\lambda)\otimes V(\mu)\longrightarrow V(\xi)$.
A \emph{shape} is a polygonal path from $0$ to $\mu$ with dominant
increments. A \emph{schedule} starts at $\lambda$, ends in the Weyl
orbit of $\xi$, and follows Weyl translates of these increments,
changing direction by suitably oriented root reflections. At a
reflection, the difference between the outgoing and incoming Weyl
translates of the shape value is a scalar multiple of the reflecting
root. This scalar is the switch coefficient; a schedule is integral
when every such coefficient is an integer.

The first construction produces a schedule whose switch roots form a
real basis. An area maximum and a torus-orbit polytope argument also
make its switches limits of Bruhat covers. A real root basis need not
generate the full root lattice, and this is where integrality can fail.
In type $D$, rank reduction leaves only one obstruction: the signed
root graph has two unbalanced components, each with an odd coordinate
sum in its assigned label difference. Joining the components introduces
one relation among the roots. Motion along half its primitive integer
generator reaches integral switch coefficients without changing the
endpoint.

Reaching those coefficients arithmetically does not yet give a path:
the switch times and signs must remain admissible throughout the motion.
Rank-two reorderings continue any locally feasible join through
simultaneous switches. A strict spectral multiplicity count, compared
with the evolution of signed coordinate comparisons, forces such a
join even at exact ties. Together these arguments give an integral
schedule whenever the moving weight has at most one zero coordinate.

A \emph{fundamental column} is a straight path segment whose increment
is a fundamental weight. The second construction converts integral
schedules for a sequence of such columns into tensor intertwiners. Minuscule columns have no
effective interior switches. An ordinary exterior column can change
only at its midpoint. A common-wall invariant covector, built from
coordinate determinants and skew pairing forms, has a nonzero
coefficient at a suitable flag. A flag-composition argument then
places the product of these coefficients on the Cartan summand for
$\mu$. This supplies the tensor occurrence without strengthening the
switch condition to a saturated path chain.

Finally, a preprojective-module construction removes a common terminal
zero block. Its main linear-algebra ingredient is a filtered
intersection estimate for the socle of a nilpotent operator. Doubling
Jordan rows with opposite alternating colors allows the top and socle
bounds to hold simultaneously at new orthogonal leaves. The doubling
is used only to obtain a smaller-rank real moment solution; the
integral schedule is then constructed at the original scale and padded
back to the original columns. These three constructions separate the
real geometry, the root-lattice obstruction, and the recovery of tensor
coefficients.

Section~\ref{sec:real-schedules} develops the real schedules and their
basis property. Sections~\ref{sec:lattice-continuation}
and~\ref{sec:moment-knots} prove their integrality.
Section~\ref{sec:columns} constructs the column intertwiners, and
Section~\ref{sec:zero-tail-completion} proves the zero-tail reduction
and completes Theorem~\ref{thm:main}.
